\diarytitle{0044}{連立線形微分方程式の一般解（相異なる実固有値の場合）}

この連立方程式は $\dfrac{d}{dt}\begin{pmatrix} x \\ y \end{pmatrix} = A \begin{pmatrix} x \\ y \end{pmatrix}$ の形であり、係数行列 $A$ の固有値・固有ベクトルを求め、それらを用いて一般解を構成する。

係数行列は $A = \begin{pmatrix} 1 & 2 \\ 2 & 1 \end{pmatrix}$。特性方程式は
\[
\det(A - \lambda I) = (1-\lambda)^{2} - 4 = \lambda^{2} - 2\lambda - 3 = (\lambda - 3)(\lambda + 1) = 0
\]
より固有値は $\lambda_{1} = 3,\ \lambda_{2} = -1$。

\noindent
$\lambda_{1} = 3$ に対して $(1-3)x + 2y = 0 \Rightarrow x = y$ より固有ベクトル $\bm{v}_{1} = \begin{pmatrix} 1 \\ 1 \end{pmatrix}$。

\noindent
$\lambda_{2} = -1$ に対して $(1+1)x + 2y = 0 \Rightarrow x = -y$ より固有ベクトル $\bm{v}_{2} = \begin{pmatrix} 1 \\ -1 \end{pmatrix}$。

\noindent
相異なる実固有値なので、一般解は $\bm{v}_{1}e^{\lambda_1 t}$ と $\bm{v}_{2}e^{\lambda_2 t}$ の線形結合として
\[
\boxed{\;
x(t) = C_{1} e^{3t} + C_{2} e^{-t}, \qquad
y(t) = C_{1} e^{3t} - C_{2} e^{-t}
\;}
\qquad (C_{1}, C_{2} \text{ は任意定数})
\]

（検算: $x' = 3C_{1}e^{3t} - C_{2}e^{-t}$ に対し $x + 2y = C_{1}e^{3t}+C_{2}e^{-t}+2C_{1}e^{3t}-2C_{2}e^{-t} = 3C_{1}e^{3t}-C_{2}e^{-t} = x'$。また $y' = 3C_{1}e^{3t}+C_{2}e^{-t}$ に対し $2x+y = 2C_{1}e^{3t}+2C_{2}e^{-t}+C_{1}e^{3t}-C_{2}e^{-t} = 3C_{1}e^{3t}+C_{2}e^{-t} = y'$。いずれも一致するので正しい。）
